\chapter{The \emph{Shijing} (Book of Songs)}

\section{What the \emph{Shijing} is}

The \zht{詩經} \emph{Shījīng} (\emph{Classic of Poetry}, \emph{Book of Songs}, \emph{Book of
Odes}) is the oldest anthology of Chinese poetry: 305 poems, most composed between about
1000 and 600 BCE in the Zhou states of the Yellow River plain. In its own time it was simply
\zht{詩} ``the Songs''; it became a \zht{經} ``classic'' when Han-dynasty scholars made it one
of the Confucian canon.

\begin{important}[Key facts for identifications]
\begin{itemize}
  \item 305 poems (often rounded to ``the three hundred''), anonymous.
  \item Predominantly \textbf{four-character lines} (\zht{四言}), grouped in stanzas, with
        heavy repetition between stanzas.
  \item Three sections: \emph{Feng} \zht{風} (Airs of the States), \emph{Ya} \zht{雅}
        (Elegantiae / Court Songs, divided into Lesser and Greater), \emph{Song} \zht{頌}
        (Hymns / Eulogies).
  \item Tradition says Confucius selected the 305 from some 3000 songs; modern scholars
        doubt this, but Confucius certainly taught the Songs (\emph{Analects}).
  \item The standard surviving text with commentary is the \emph{Mao} version
        (\zht{毛詩}), whose ``Great Preface'' (\zht{大序}) gives the earliest theory of
        poetry.
\end{itemize}
\end{important}

\section{The three sections}

\begin{center}
\begin{tabular}{llp{0.5\textwidth}}
\toprule
Section & Number & Character \\
\midrule
\emph{Guofeng} \zht{國風} ``Airs of the States'' & 160 & Songs from fifteen regional states;
  courtship, marriage, farming, soldiers' laments; closest to ``folk song'' (though
  polished by court musicians) \\
\emph{Xiaoya} \zht{小雅} ``Lesser Elegantiae'' & 74 & Court songs: banquets, campaigns,
  complaints of officials \\
\emph{Daya} \zht{大雅} ``Greater Elegantiae'' & 31 & Dynastic narratives, the founding of
  Zhou, King Wen \\
\emph{Song} \zht{頌} ``Hymns'' & 40 & Sacrificial hymns of Zhou, Lu and Shang, for ancestral
  temple ritual, often with dance \\
\bottomrule
\end{tabular}
\end{center}

\begin{remark}
The three names may originally have been musical categories (regional tunes, court
``correct'' music, temple music with dance) rather than thematic ones. The music is lost;
only the words survive.
\end{remark}

\section{The Great Preface: what poetry is}

The Great Preface to the Mao edition (probably Han) contains the most quoted definition of
poetry in Chinese:

\begin{center}
\zht{詩者，志之所之也。在心為志，發言為詩。}
\end{center}

My rendering: ``Poetry is where the intent goes. In the heart it is intent (\zht{志}
\emph{zhì}); expressed in words it is poetry.'' The Preface continues that when words are not
enough, one sighs; when sighing is not enough, one sings; when singing is not enough, one's
hands and feet begin to dance.

\begin{important}[\emph{Shi yan zhi} \zht{詩言志}]
``Poetry articulates what is on the mind intently.'' This formula (already in the
\emph{Book of Documents}) is the foundation of the Chinese view that a poem is a
\emph{sincere expression} of the poet's inner state, prompted by the world. Contrast the
Aristotelian Western starting point of poetry as \emph{mimesis} (imitation, fiction). One
consequence: Chinese readers have usually read poems \emph{biographically}, as true
statements by a real person on a real occasion.
\end{important}

The Preface also gives the \textbf{six principles} (\zht{六義}): \emph{feng}, \emph{ya},
\emph{song} (the three genres), and \emph{fu}, \emph{bi}, \emph{xing} (the three modes of
presentation).

\section{\emph{Fu}, \emph{bi}, \emph{xing}}

\begin{definition}[Fu, bi, xing]
\begin{itemize}
  \item \zht{賦} \textbf{\emph{fù}} ``exposition'': direct statement or description of a
        thing or event.
  \item \zht{比} \textbf{\emph{bǐ}} ``comparison'': an explicit or clearly intended
        likeness (simile / metaphor): A is like B.
  \item \zht{興} \textbf{\emph{xìng}} ``stirring, evocation'': a natural image placed at
        the head of a stanza that \emph{arouses} the feeling of what follows, without the
        relation being spelled out.
\end{itemize}
\end{definition}

\begin{center}
\begin{tikzpicture}[box/.style={draw, thick, rounded corners, align=center,
    minimum width=3cm, minimum height=1cm, font=\small}]
  \node[box, fill=gray!10] (fu) at (0,0) {\zht{賦} \emph{fu}\\ say it directly};
  \node[box, fill=orange!15] (bi) at (5.2,0) {\zht{比} \emph{bi}\\ A is like B};
  \node[box, fill=green!15] (xing) at (10.4,0) {\zht{興} \emph{xing}\\ image, then feeling};
  \node[font=\scriptsize, align=center] at (0,-1.2) {``We gathered ferns\\ on the hill''};
  \node[font=\scriptsize, align=center] at (5.2,-1.2) {``Big rat, big rat,\\ don't eat our millet''\\ (rat = the lord)};
  \node[font=\scriptsize, align=center] at (10.4,-1.2) {``Ospreys cry on the islet''\\ $\leadsto$ ``the gentleman\\ seeks a fair maiden''};
  \draw[-{Stealth}, thick] (fu) -- node[above, font=\scriptsize] {explicitness} (bi);
  \draw[{Stealth}-, thick] (bi) -- node[above, font=\scriptsize] {decreases} (xing);
\end{tikzpicture}
\end{center}

\begin{remark}
\emph{Xing} is the distinctive one and the hardest to pin down. Is the link between the
opening image and the human situation a comparison (birds paired $\Rightarrow$ people
paired)? An association of the season (spring $\Rightarrow$ weddings)? Or originally just a
conventional musical opening with a rhyme? Scholars disagree; for exams, say that
\emph{xing} \emph{evokes} rather than \emph{states} the comparison.
\end{remark}

\section{Reading ``Guan ju'', the first poem}

The first poem of the anthology, ``Guan ju'' \zht{關雎} (Mao 1), is named, like most
\emph{Shijing} poems, after its opening words.

\begin{center}
\begin{tabular}{p{0.22\textwidth}p{0.28\textwidth}p{0.4\textwidth}}
\toprule
Original & Pinyin & My translation \\
\midrule
\zht{關關雎鳩，} & guān guān jū jiū & ``Guan, guan'' cry the ospreys \\
\zht{在河之洲。} & zài hé zhī zhōu & on the islet in the river. \\
\zht{窈窕淑女，} & yǎotiǎo shū nǚ & The gentle, lovely maiden: \\
\zht{君子好逑。} & jūnzǐ hǎo qiú & a fit match for our lord. \\
\midrule
\zht{參差荇菜，} & cēncī xìng cài & Long and short, the floating-heart, \\
\zht{左右流之。} & zuǒ yòu liú zhī & to left and right we seek it. \\
\zht{窈窕淑女，} & yǎotiǎo shū nǚ & The gentle, lovely maiden: \\
\zht{寤寐求之。} & wù mèi qiú zhī & waking and sleeping he sought her. \\
\bottomrule
\end{tabular}
\end{center}

The poem continues: he seeks her and cannot find her, tosses and turns all night; then the
last stanzas gather the water plant ``to left and right'' and welcome her with zithers,
lutes, bells and drums.

\subsection{Analysis}

\begin{itemize}
  \item \textbf{Form.} Four-character lines; rhymes \emph{jiū / zhōu / qiú} (line 1, 2, 4).
        Stanzas repeat with small substitutions (\zht{流} ``seek'' $\to$ \zht{采}
        ``pick'' $\to$ \zht{芼} ``select''), a typical incremental repetition of song.
  \item \textbf{Sound.} \zht{關關} is onomatopoeia (the bird call). \zht{窈窕}
        \emph{yǎotiǎo} and \zht{參差} \emph{cēncī} are \emph{binomes} (\zht{聯綿詞}),
        two-syllable words linked by rhyme or alliteration, a hallmark of \emph{Shijing}
        diction.
  \item \textbf{\emph{Xing}.} The ospreys (believed to mate for life, and to keep
        apart in public) open the poem and evoke the ideal couple. The water plant being
        gathered (a task for the women preparing ancestral offerings) evokes the
        courtship and the bride's future ritual role.
  \item \textbf{Interpretation.} Read plainly, it is a courtship/wedding song. The Mao
        commentary read it as praising the virtue of the queen consort of King Wen, who
        sought worthy ladies for the harem without jealousy. That allegorical reading
        made a love song suitable to open a Confucian classic.
\end{itemize}

\begin{aside}
Confucius, \emph{Analects} 3.20: ```Guan ju' is joyful without excess, mournful without
injury.'' This became the ideal of emotional balance in poetry: feeling expressed, but
regulated.
\end{aside}

\section{More poems from the \emph{Airs}}

\subsection{``The Peach Tree is Young'' \zht{桃夭} (Mao 6)}

\begin{center}
\begin{tabular}{p{0.3\textwidth}p{0.55\textwidth}}
\toprule
\zht{桃之夭夭，} & The peach tree, young and tender, \\
\zht{灼灼其華。} & blazing, blazing are its flowers. \\
\zht{之子于歸，} & This girl is going to her new home: \\
\zht{宜其室家。} & she will be good for her chamber and her house. \\
\bottomrule
\end{tabular}
\end{center}

The second and third stanzas move from flowers to \emph{fruit} (\zht{有蕡其實}) to
\emph{thick leaves} (\zht{其葉蓁蓁}): a wedding song whose \emph{xing} images grow with
the bride's expected life (flowering, fertility, a flourishing household).
\zht{于歸} ``going home'' is the standard term for a woman marrying: her ``home'' is
her husband's.

\subsection{Complaint and satire}

Not all the \emph{Airs} are happy. ``Big Rat'' (\zht{碩鼠}, Mao 113) addresses a rat
eating the millet, which every reader since the Han has taken as a \emph{bi} for a greedy
ruler, ending with the farmers' wish to leave for a ``happy land''. Soldiers' songs (e.g.
``We Gather Ferns'' \zht{采薇}, Lesser \emph{Ya}) end with the famous lines on returning
home:

\begin{center}
\begin{tabular}{p{0.3\textwidth}p{0.55\textwidth}}
\toprule
\zht{昔我往矣，} & Long ago, when I set out, \\
\zht{楊柳依依。} & the willows were lush and clinging. \\
\zht{今我來思，} & Now as I come back, \\
\zht{雨雪霏霏。} & the snow falls thick and heavy. \\
\bottomrule
\end{tabular}
\end{center}

Here the willow, already linked with parting, makes its first appearance; the contrast of
spring departure and winter return compresses years of campaigning into two images.

\section{Formal features to recognise}

\begin{center}
\begin{tabular}{lp{0.65\textwidth}}
\toprule
Feature & Description \\
\midrule
Four-character line & The basic rhythm: two beats of two syllables \\
Stanzaic repetition & Stanzas differ by one or two words (incremental repetition) \\
Reduplicatives & \zht{關關}, \zht{灼灼}, \zht{依依}, \zht{霏霏}: doubled syllables for
  sound or vividness \\
Binomes & \zht{窈窕}, \zht{參差}: two-syllable rhyming / alliterating words \\
Particles & \zht{兮}, \zht{之}, \zht{矣}, \zht{思} fill out the rhythm \\
Opening \emph{xing} image & Nature first, then the human situation \\
\bottomrule
\end{tabular}
\end{center}

\section*{Summary}
\begin{itemize}
  \item The \emph{Shijing}: 305 anonymous poems, c.\ 1000--600 BCE, mostly four-character
        lines; sections \emph{feng}, \emph{ya}, \emph{song}.
  \item The Great Preface: ``poetry articulates intent'' (\zht{詩言志}), poetry as
        sincere expression; the six principles.
  \item \emph{Fu} (exposition), \emph{bi} (comparison), \emph{xing} (evocative opening
        image).
  \item ``Guan ju'': osprey \emph{xing}, binomes, wedding song read allegorically by the
        Mao commentary.
  \item The Confucian reading turned folk and court songs into moral instruction; this
        habit of allegorical and biographical reading persists through the tradition.
\end{itemize}
